\section{Graphs}
% Covers lecture notes chapter 6 (Sep 29 build of the notes).

% ---- helper macros for this chapter (prefix \gr) ----
% \grBase: the running example graph G (vertices a..f) in light gray, nodes named a..f.
\newcommand{\grBase}{%
  \node[vert] (a) at (0,0) {$a$};
  \node[vert] (b) at (1.6,0) {$b$};
  \node[vert] (c) at (0.8,-1.1) {$c$};
  \node[vert] (d) at (0.8,-2.3) {$d$};
  \node[vert] (e) at (2.75,-0.55) {$e$};
  \node[vert] (f) at (2.75,-1.75) {$f$};
  \draw[gray] (a)--(b) (b)--(c) (c)--(a) (c)--(d) (e)--(f);}
% \grStep{node A}{node B}{label}: step number on the edge between A and B
\newcommand{\grStep}[3]{\node[font=\scriptsize,fill=white,inner sep=1pt] at ($(#1)!0.5!(#2)$) {#3};}

\begin{notebox}
\textbf{Is this on Test 1?} Graphs were lectured the week before the test. Ask your TA whether this chapter is on Test~1 before you spend your last study night on it. Either way it comes back later in the course, so the time isn't wasted.
\end{notebox}

A graph is the formal version of a dots-and-lines picture: a network of friends, cities joined by roads, web pages joined by links. The math keeps only two facts: which dots exist, and which pairs of dots are joined. Everything else about the picture is decoration.

% ------------------------------------------------------------------
\subsection{What a graph is}

\begin{notebox}
\textbf{Graph.} A graph is an ordered pair $G=(V,E)$. $V$ is a set (the \emph{vertices}, or nodes). $E$ is a set of two-member sets of vertices (the \emph{edges}).
\end{notebox}

Every piece of that definition is doing a job.
\begin{itemize}
\item The pair $(V,E)$ is ordered because the two parts play different roles.
\item $V$ and $E$ are sets, so vertices and edges have no order and no repeats. You can't have two separate edges between the same two vertices.
\item Each edge is a set, so an edge has no direction: $\{v,w\}=\{w,v\}$.
\item Each edge has exactly \emph{two} members. An edge from $v$ to itself would be $\{v,v\}=\{v\}$, a one-member set, so it isn't allowed. No loops in a plain graph.
\end{itemize}

Position, edge length, curves and crossings are all artifacts of the drawing. Two drawings with the same vertices and the same joined pairs are the same graph.

\begin{figure}[H]
\centering
\begin{tikzpicture}
  \begin{scope}
    \node[vert] (1) at (0,0) {1}; \node[vert] (2) at (2,0) {2};
    \node[vert] (3) at (2,-2) {3}; \node[vert] (4) at (0,-2) {4};
    \draw (1)--(2) (2)--(3) (3)--(4) (4)--(1) (1)--(3);
  \end{scope}
  \begin{scope}[xshift=5.2cm]
    \node[vert] (1) at (1,0) {1}; \node[vert] (3) at (1,-2) {3};
    \node[vert] (2) at (0,-1) {2}; \node[vert] (4) at (3,-1) {4};
    \draw (1)--(2) (2)--(3) (3)--(4) (4)--(1) (1)--(3);
  \end{scope}
  \node[font=\footnotesize,align=left,anchor=west] at (9.6,-1) {$V=\{1,2,3,4\}$\\$E=\{\{1,2\},\{2,3\},\{3,4\},$\\$\phantom{E=\{}\{4,1\},\{1,3\}\}$};
\end{tikzpicture}
\caption{Two drawings, one graph. Same vertices, same five joined pairs.}
\end{figure}

\textbf{Worked example.} Which of these are graphs?
\begin{itemize}
\item $(\{1,2,3\},\{\{1,2\},\{2,3\}\})$. Yes.
\item $(\{1,2\},\{\{1,2\},\{2,1\}\})$. Yes, and it has one edge, not two. $\{1,2\}$ and $\{2,1\}$ are the same set, so $E$ has one member.
\item $(\{1,2,3\},\{\{1,1\},\{2,3\}\})$. No. $\{1,1\}=\{1\}$ isn't a two-member set.
\item $(\{1,2\},\{(1,2)\})$. Not a graph. $(1,2)$ is an ordered pair, not a set; this is a \emph{directed} graph (see below).
\end{itemize}

% ------------------------------------------------------------------
\subsection{Adjacency and degree}

\begin{notebox}
If $\{v,w\}\in E$, the edge \emph{joins} $v$ and $w$, and $v$ and $w$ are \emph{adjacent}; each is a \emph{neighbor} of the other.

\textbf{Degree of a vertex} $v$: the number of neighbors of $v$ (the same as the number of edges that include $v$).

\textbf{Degree of the graph} $G$: the \emph{maximum} degree of any vertex in $G$.
\end{notebox}

A vertex is never adjacent to itself in a plain graph, because there are no loops.

Here's the running example for the rest of the chapter:
\[V=\{a,b,c,d,e,f\},\qquad E=\{\{a,b\},\{b,c\},\{c,a\},\{c,d\},\{e,f\}\}.\]

\begin{figure}[H]
\centering
\begin{tikzpicture}[x=1.35cm,y=1.35cm,dg/.style={font=\scriptsize,inner sep=1pt}]
  \node[vert] (a) at (0,0) {$a$};
  \node[vert] (b) at (1.6,0) {$b$};
  \node[vert] (c) at (0.8,-1) {$c$};
  \node[vert] (d) at (0.8,-2.1) {$d$};
  \node[vert] (e) at (3.3,-0.5) {$e$};
  \node[vert] (f) at (3.3,-1.7) {$f$};
  \draw (a)--(b) (b)--(c) (c)--(a) (c)--(d) (e)--(f);
  \node[dg,left=3pt of a] {deg 2};
  \node[dg,right=3pt of b] {deg 2};
  \node[dg,left=3pt of c] {deg 3};
  \node[dg,left=3pt of d] {deg 1};
  \node[dg,right=3pt of e] {deg 1};
  \node[dg,right=3pt of f] {deg 1};
\end{tikzpicture}
\caption{The running example $G$. Degrees $2,2,3,1,1,1$, so the degree of $G$ is $3$: the largest, not the total.}
\end{figure}

\textbf{Worked example.} The neighbors of $c$ are $a$, $b$ and $d$, so $\deg(c)=3$. $a$ and $d$ aren't adjacent; there's no edge $\{a,d\}$. The degree of $G$ is $\max(2,2,3,1,1,1)=3$.

\textbf{Slip.} ``Degree of the graph'' is the biggest vertex degree, not the sum. The sum is still a handy check, though it isn't something the notes define: each edge adds $1$ to the degree of each of its two ends, so the vertex degrees always add up to twice the number of edges. Here $2+2+3+1+1+1=10=2\cdot 5$. If your degrees don't add to an even number, you've miscounted.

% ------------------------------------------------------------------
\subsection{Directed graphs}

Sometimes direction matters. A link from page 1 to page 2 isn't a link from page 2 to page 1. So swap the two-member sets for ordered pairs.

\begin{notebox}
\textbf{Directed graph.} $G=(V,E)$ where $V$ is a set and $E$ is a set of \emph{ordered pairs} of vertices.
\begin{itemize}
\item The edge $(v,w)$ goes \emph{from} $v$ \emph{to} $w$. $v$ is the \emph{tail}, $w$ is the \emph{head}. (The arrowhead sits on the head.)
\item $w$ is a \emph{direct successor} (child) of $v$; $v$ is a \emph{direct predecessor} (parent) of $w$.
\item Loops $(v,v)$ are allowed, and $(v,w)$ and $(w,v)$ are different edges.
\item \textbf{Indegree} of $v$: the number of direct predecessors (edges pointing into $v$). \textbf{Outdegree}: the number of direct successors (edges pointing out of $v$).
\end{itemize}
\end{notebox}

\begin{figure}[H]
\centering
\begin{tikzpicture}[dg/.style={font=\scriptsize,inner sep=1pt,align=center}]
  \node[vert] (1) at (0,0) {1};
  \node[vert] (2) at (3,0) {2};
  \node[vert] (3) at (3,-2.2) {3};
  \node[vert] (4) at (0,-2.2) {4};
  \draw[->] (1) -- (2);
  \draw[->] (2) -- (3);
  \draw[->] (3) -- (1);
  \draw[->] (1) -- (4);
  \draw[->] (4) -- (3);
  \draw[->] (3) to[out=-30,in=-90,looseness=6] (3);
  \node[dg,anchor=south] at (1.north) {in 1, out 2};
  \node[dg,anchor=south] at (2.north) {in 1, out 1};
  \node[dg,anchor=west] at ($(3.east)+(0.75,0.25)$) {in 3, out 2};
  \node[dg,anchor=north] at (4.south) {in 1, out 1};
  \node[font=\footnotesize,align=left,anchor=west] at (6.0,-1.1) {$E=\{(1,2),(2,3),(3,1),$\\$\phantom{E=\{}(1,4),(4,3),(3,3)\}$\\[3pt]edge $(2,3)$: tail $2$, head $3$\\the loop $(3,3)$ counts once in\\and once out for vertex $3$};
\end{tikzpicture}
\caption{A directed graph with each vertex's indegree and outdegree. Both columns add up to $6$, the number of edges, since every edge has one tail and one head.}
\end{figure}

\textbf{Worked example.} Vertex $3$ has direct predecessors $2$, $4$ and $3$ itself (from the loop), so its indegree is $3$. Its direct successors are $1$ and $3$, so its outdegree is $2$. Is $(2,1)$ an edge? No. $(1,2)$ is, and in a directed graph those are different edges.

\textbf{Slip.} If your copy of the notes says the head is the starting end, it's an older build. The current notes, and the usual convention, say the edge $(v,w)$ goes from tail $v$ to head $w$.

\subsubsection{Weighted graphs}

When edges need a number attached (a distance, a travel time, a bandwidth), you add a weight function. A weighted directed graph is a triple $(V,E,w)$ where $w:E\to N$ assigns each edge a weight from some set $N$. Usually $N=\mathbb{R}^{+}$, the positive \emph{real} numbers. (If your copy calls $\mathbb{R}^{+}$ the positive integers, that's a slip.) Weighted graphs barely come up in this course; just know what the triple means.

% ------------------------------------------------------------------
\subsection{Walks, trails and paths}

Back to plain graphs. A walk is a route you could trace with a pencil without lifting it, moving only along edges. Trails and paths add restrictions.

\begin{notebox}
\textbf{Walk} from $v$ to $w$: a sequence of edges joining a sequence of vertices $v_0,v_1,\dots,v_n$ with $v_0=v$ and $v_n=w$, where each edge joins $v_{i-1}$ and $v_i$. We usually write the vertex sequence, like $(a,b,c)$.

\textbf{Trail:} a walk with no repeated edge.

\textbf{Path:} a walk with no repeated vertex.
\end{notebox}

Using an edge twice means visiting one of its ends twice, so a repeated edge always brings a repeated vertex with it. The reverse fails: $(d,c,a,b,c)$ below repeats $c$ without repeating an edge. So every path is a trail, and every trail is a walk. Each word is a stricter version of the one before.

Two edge cases. An edge counts as repeated even if you cross it in the opposite direction, because $\{a,b\}=\{b,a\}$. And $(v)$ on its own is the \emph{trivial path}: zero edges, a path from $v$ to itself.

\begin{figure}[H]
\centering
\begin{tikzpicture}[x=1.0cm,y=1.0cm]
  \newcommand{\grPanel}[4]{%
    \begin{scope}[xshift=#1]
      \grBase
      #3
      \node[font=\small,anchor=south] at (1.35,0.35) {#2};
      \node[font=\footnotesize,align=center,anchor=north,text width=3.5cm] at (1.35,-2.75) {#4};
    \end{scope}}
  \grPanel{0cm}{walk, not trail}{\draw[hl] (a)--(b)--(c)--(a);
      \node[font=\scriptsize,fill=white,inner sep=1pt,above] at ($(a)!0.5!(b)$) {1, 4};
      \grStep{b}{c}{2}\grStep{c}{a}{3}}{$(a,b,c,a,b)$\\edge $\{a,b\}$ used twice}
  \grPanel{4.05cm}{trail, not path}{\draw[hl] (d)--(c)--(a)--(b)--(c);
      \grStep{d}{c}{1}\grStep{c}{a}{2}\grStep{a}{b}{3}\grStep{b}{c}{4}}{$(d,c,a,b,c)$\\no edge repeats; vertex $c$ does}
  \grPanel{8.1cm}{path}{\draw[hl] (d)--(c)--(b)--(a);
      \grStep{d}{c}{1}\grStep{c}{b}{2}\grStep{b}{a}{3}}{$(d,c,b,a)$\\no repeated vertex}
  \grPanel{12.15cm}{cycle}{\draw[hl] (a)--(b)--(c)--(a);
      \grStep{a}{b}{1}\grStep{b}{c}{2}\grStep{c}{a}{3}}{$(a,b,c,a)$: non-empty trail,\\only first $=$ last repeats}
\end{tikzpicture}
\caption{The running example four times. Bold edges are the route; the numbers give the order. In the first panel, edge $\{a,b\}$ is step 1 and step 4.}
\end{figure}

To classify a vertex sequence, ask these questions in order and stop at the first ``no'':
\begin{enumerate}
\item Is every consecutive pair joined by an edge? If not, it isn't a walk at all.
\item Is any edge used twice (in either direction)? If so, it's a walk but not a trail.
\item Is any vertex visited twice? If so, it's a trail but not a path. If not, it's a path.
\end{enumerate}

\textbf{Worked example} (running example $G$).
\begin{center}
\begin{tabularx}{\linewidth}{@{}llL@{}}
\toprule
Sequence & Verdict & Reason\\
\midrule
$(a,d)$ & not a walk & no edge $\{a,d\}$\\
$(a,b,c,a,b)$ & walk, not trail & edge $\{a,b\}$ repeats\\
$(d,c,a,b,c)$ & trail, not path & no edge repeats, but vertex $c$ does\\
$(d,c,b,a)$ & path & no repeated vertex\\
$(e,f)$ & path & one edge still counts\\
$(c)$ & trivial path & zero edges\\
\bottomrule
\end{tabularx}
\end{center}

\textbf{Slip.} At least one build of the notes has the sentence ``since every walk is a path'' in the connectivity section. It's backwards. Every \emph{path} is a walk; $(a,b,c,a,b)$ is a walk that isn't a path.

% ------------------------------------------------------------------
\subsection{Connectivity}

\begin{notebox}
\textbf{Connected vertices.} $v$ is connected to $w$ if and only if there is a path from $v$ to $w$. (A walk works just as well: if there's a walk, there's a path.)

\textbf{Connected graph.} Every pair of vertices is connected.

\textbf{Connected component.} A \emph{maximal} set of vertices that are all connected to each other. ``Maximal'' means you can't add another vertex and keep it connected.
\end{notebox}

\begin{figure}[H]
\centering
\begin{tikzpicture}[x=1.35cm,y=1.35cm]
  \node[vert] (a) at (0,0) {$a$};
  \node[vert] (b) at (1.6,0) {$b$};
  \node[vert] (c) at (0.8,-1) {$c$};
  \node[vert] (d) at (0.8,-2.1) {$d$};
  \node[vert] (e) at (3.3,-0.5) {$e$};
  \node[vert] (f) at (3.3,-1.7) {$f$};
  \draw (a)--(b) (b)--(c) (c)--(a) (c)--(d) (e)--(f);
  \node[draw,dashed,gray,rounded corners=6pt,fit=(a)(b)(c)(d),inner sep=7pt] (K1) {};
  \node[draw,dashed,gray,rounded corners=6pt,fit=(e)(f),inner sep=7pt] (K2) {};
  \node[font=\scriptsize,anchor=south] at (K1.north) {component $\{a,b,c,d\}$};
  \node[font=\scriptsize,anchor=south] at (K2.north) {component $\{e,f\}$};
\end{tikzpicture}
\caption{$G$ has two connected components, so $G$ is not connected: there's no path from $a$ to $e$.}
\end{figure}

\textbf{Worked example.}
\begin{itemize}
\item $d$ is connected to $b$: $(d,c,b)$ is a path.
\item $a$ is not connected to $f$, so $G$ isn't a connected graph.
\item $\{a,b\}$ is not a component. Its vertices are connected, but you can add $c$ and stay connected, so it isn't maximal.
\item $\{a,e\}$ is not a component either; $a$ and $e$ aren't even connected.
\item Every vertex is connected to itself by the trivial path. So an isolated vertex with no edges is a component on its own.
\end{itemize}

Side note: ``is connected to'' is reflexive, symmetric and transitive. A relation with all three is an \emph{equivalence relation}, and components are exactly its groups. In a directed graph you also need to say which kind: \emph{strongly} connected follows the arrows, \emph{weakly} connected ignores their direction.

% ------------------------------------------------------------------
\subsection{Closed walks, circuits and cycles}

\begin{notebox}
\begin{itemize}
\item \textbf{Closed walk:} a walk whose first and last vertices are the same.
\item \textbf{Circuit:} a non-empty trail whose first and last vertices are the same.
\item \textbf{Cycle:} a non-empty trail where the only repeated vertex is first $=$ last.
\end{itemize}
``Non-empty'' means at least one edge.
\end{notebox}

Why isn't a cycle just ``a closed path''? Because a path can't repeat any vertex, and a closed route repeats its starting vertex at the end. So the definition starts from a trail and bans every other repeat.

In a plain graph the shortest cycle has three edges. One edge would need a loop, which isn't allowed. Two edges would look like $(v,w,v)$, which uses $\{v,w\}$ twice, so it isn't even a trail.

\textbf{Worked example} (running example $G$).
\begin{itemize}
\item $(a,b,c,a)$ is a cycle, and also a circuit and a closed walk.
\item $(c,a,c)$ is a closed walk only. It repeats the edge $\{c,a\}$.
\item $(c)$ is a closed walk (trivially), but not a circuit or a cycle: it has no edges.
\item $(c,a,b,c,d)$ is not closed at all, since it starts at $c$ and ends at $d$.
\end{itemize}
For a circuit that isn't a cycle you need a vertex visited twice in the middle without reusing an edge, which takes two triangles sharing a vertex. Add $\{c,g\}$, $\{g,h\}$ and $\{h,c\}$ to $G$; then $(a,b,c,g,h,c,a)$ is a circuit but not a cycle, because $c$ repeats in the middle.

% ------------------------------------------------------------------
\subsection{Trees}

\begin{notebox}
\textbf{Tree.} A tree is a connected graph with no cycles.

\textbf{Theorem (notes 6.6.1).} A graph is a tree if and only if every pair of vertices has exactly one path between them.
\end{notebox}

That's the whole official definition. Roots, parents and leaves aren't in it; they come from choosing a root afterwards. The theorem gives you a second test that's often quicker to apply: if you can find two different paths between some pair of vertices, there's a cycle hiding there.

\begin{figure}[H]
\centering
\begin{tikzpicture}
  \begin{scope}
    \node[vert] (1) at (0,0) {1}; \node[vert] (2) at (-0.9,-1.2) {2}; \node[vert] (3) at (0.9,-1.2) {3};
    \node[vert] (4) at (-0.9,-2.4) {4}; \node[vert] (5) at (0.9,-2.4) {5};
    \draw (1)--(2) (1)--(3) (2)--(4) (3)--(5);
    \node[font=\footnotesize,align=center,anchor=north] at (0,-2.95) {tree\\connected, no cycle};
  \end{scope}
  \begin{scope}[xshift=4.6cm]
    \node[vert] (1) at (0,0) {1}; \node[vert] (2) at (-0.9,-1.2) {2}; \node[vert] (3) at (0.9,-1.2) {3};
    \node[vert] (4) at (-0.9,-2.4) {4}; \node[vert] (5) at (0.9,-2.4) {5};
    \draw (1)--(2) (1)--(3) (2)--(4) (3)--(5);
    \draw[hl] (2)--(3);
    \draw[hl] (1)--(2) (1)--(3);
    \node[font=\footnotesize,align=center,anchor=north] at (0,-2.95) {not a tree\\cycle $(1,2,3,1)$};
  \end{scope}
  \begin{scope}[xshift=9.2cm]
    \node[vert] (1) at (0,0) {1}; \node[vert] (2) at (-0.9,-1.2) {2}; \node[vert] (3) at (0.9,-1.2) {3};
    \node[vert] (4) at (-0.9,-2.4) {4}; \node[vert] (5) at (0.9,-2.4) {5};
    \draw (1)--(2) (1)--(3) (2)--(4);
    \node[font=\footnotesize,align=center,anchor=north] at (0,-2.95) {not a tree\\$5$ is cut off};
  \end{scope}
\end{tikzpicture}
\caption{Adding one edge to a tree creates a cycle (middle). Removing one disconnects it (right).}
\end{figure}

\subsubsection{Rooted trees}

Pick one vertex as the root. Because there's exactly one path from each vertex to the root, everything else gets a precise meaning.

\begin{notebox}
\textbf{Rooted tree:} a tree with one vertex chosen as the root. Then:
\begin{itemize}
\item the \emph{parent} of $v$ is the neighbor of $v$ on the path from $v$ to the root, and $v$ is a \emph{child} of its parent;
\item a \emph{leaf} is a vertex with no children;
\item a \emph{branch} is a path from the root to a leaf (the whole route, not one fork);
\item the \emph{depth} of $v$ is the number of \textbf{edges} on the path from $v$ to the root;
\item the \emph{height} of $v$ is the number of \textbf{edges} on the longest path from $v$ that goes away from the root down to a leaf;
\item the \emph{depth of the tree} is the largest depth of any vertex (the same as the root's height, and the length of the longest branch).
\end{itemize}
\end{notebox}

\begin{figure}[H]
\centering
\begin{tikzpicture}[x=1cm,y=1.05cm]
  \foreach \l in {0,1,2,3} {
    \draw[pruned] (-0.6,-\l) -- (5.0,-\l);
    \node[font=\footnotesize,anchor=east] at (-0.7,-\l) {depth \l};
  }
  \node[vert,fill=white] (r) at (2.4,0) {$r$};
  \node[vert,fill=white] (s) at (1.2,-1) {$s$};
  \node[vert,fill=white] (t) at (3.8,-1) {$t$};
  \node[vert,fill=white] (u) at (0.4,-2) {$u$};
  \node[vert,fill=white] (v) at (2.0,-2) {$v$};
  \node[vert,fill=white] (w) at (3.8,-2) {$w$};
  \node[vert,fill=white] (x) at (2.0,-3) {$x$};
  \draw (s)--(u) (r)--(t) (t)--(w);
  \draw[hl] (r)--(s) (s)--(v) (v)--(x);
  \node[font=\scriptsize,fill=white,inner sep=1pt] at ($(r)!0.5!(s)+(-0.2,0.05)$) {1};
  \node[font=\scriptsize,fill=white,inner sep=1pt] at ($(s)!0.5!(v)+(0.22,0)$) {2};
  \node[font=\scriptsize,fill=white,inner sep=1pt] at ($(v)!0.5!(x)+(0.22,0)$) {3};
  \node[font=\footnotesize,align=left,anchor=west] at (5.6,-1.5) {root $r$, depth $0$\\leaves $u$, $x$, $w$\\depth of $x$ $=3$ (3 edges; 4 vertices)\\height of $s=2$, via $(s,v,x)$\\height of $t=1$; height of $r=3$\\depth of the tree $=3$};
\end{tikzpicture}
\caption{A rooted tree. The bold branch $(r,s,v,x)$ is the longest one. Count the edges, not the circles.}
\end{figure}

\textbf{Worked example} (the tree above).
\begin{itemize}
\item The parent of $v$ is $s$, since $s$ is the next vertex on the path $(v,s,r)$ to the root. The children of $s$ are $u$ and $v$.
\item The branches are $(r,s,u)$, $(r,s,v,x)$ and $(r,t,w)$: one per leaf.
\item depth$(v)=2$. height$(v)=1$. height$(u)=0$, like every leaf.
\item Any vertex can be the root. Re-root at $w$ and the depth of $x$ becomes $5$, along $(x,v,s,r,t,w)$. Same tree, different rooted tree.
\end{itemize}

\textbf{Slip.} Depth and height count edges. A branch with six vertices has five edges, so it gives depth $5$, not $6$. If an example in your copy of the notes counts vertices for depth, use the edge count.

\subsubsection{$n$-ary, full and complete}

\begin{notebox}
\textbf{$n$-ary tree:} every vertex has at most $n$ children (binary when $n=2$, ternary when $n=3$).

\textbf{Full $n$-ary tree:} every vertex that isn't a leaf has exactly $n$ children.

\textbf{Complete $n$-ary tree:} a full $n$-ary tree whose leaves all have the same depth.
\end{notebox}

So complete implies full, and full implies $n$-ary, but not the other way round. A single vertex counts as a binary tree that is full and complete: it has no non-leaf vertices to fail the rule, and its only leaf is at depth $0$.

\begin{figure}[H]
\centering
\begin{tikzpicture}[x=0.8cm,y=0.95cm,sv/.style={vert,minimum size=1.3em,font=\scriptsize}]
  \begin{scope}
    \node[sv] (0) at (0,0) {}; \node[sv] (1) at (-1,-1) {}; \node[sv] (2) at (1,-1) {};
    \node[sv] (3) at (-1,-2) {};
    \draw (0)--(1) (0)--(2) (1)--(3);
    \node[draw,dashed,circle,inner sep=6pt] at (1) {};
    \node[font=\footnotesize,align=center,anchor=north] at (0,-2.6) {binary, not full\\one node has 1 child};
  \end{scope}
  \begin{scope}[xshift=5.4cm]
    \node[sv] (0) at (0,0) {}; \node[sv] (1) at (-1,-1) {}; \node[sv] (2) at (1,-1) {};
    \node[sv] (3) at (-1.6,-2) {}; \node[sv] (4) at (-0.4,-2) {};
    \draw (0)--(1) (0)--(2) (1)--(3) (1)--(4);
    \node[font=\footnotesize,align=center,anchor=north] at (0,-2.6) {full, not complete\\leaves at depths 2, 2, 1};
  \end{scope}
  \begin{scope}[xshift=10.8cm]
    \node[sv] (0) at (0,0) {}; \node[sv] (1) at (-1,-1) {}; \node[sv] (2) at (1,-1) {};
    \node[sv] (3) at (-1.6,-2) {}; \node[sv] (4) at (-0.4,-2) {};
    \node[sv] (5) at (0.4,-2) {}; \node[sv] (6) at (1.6,-2) {};
    \draw (0)--(1) (0)--(2) (1)--(3) (1)--(4) (2)--(5) (2)--(6);
    \node[font=\footnotesize,align=center,anchor=north] at (0,-2.6) {complete\\full, all leaves at depth 2};
  \end{scope}
\end{tikzpicture}
\caption{Three binary trees, each rooted at the top vertex. The dashed ring marks the vertex that breaks ``full.''}
\end{figure}

\textbf{Slip.} The full $n$-ary definition is about non-leaf vertices. If a sentence in your copy says ``every node has exactly $n$ children,'' it means every \emph{non-leaf} node; leaves have none by definition.

% ------------------------------------------------------------------
\needspace{8\baselineskip}
\subsection{Practice}

Problems 2--4 use the graph $H$ with $V=\{p,q,r,s,t,u\}$ and $E=\{\{p,q\},\{p,r\},\{p,s\},\{q,r\},\{s,t\}\}$. Draw it first.

\begin{enumerate}
\item Which of these are graphs?\\
(a) $(\{1,2,3\},\{\{1,2\},\{1,3\}\})$\quad (b) $(\{1,2\},\{\{1,2\},\{2,1\}\})$\quad (c) $(\{1,2\},\{\{1,3\}\})$\quad (d) $(\{1,2,3\},\{\{2,2\},\{1,3\}\})$
\item In $H$: give the degree of every vertex and the degree of $H$. Check your work with the degree sum.
\item In $H$, classify each sequence as not a walk, walk but not trail, trail but not path, or path. Then say which are closed walks, circuits or cycles.
\begin{center}
(a) $(q,p,s,t)$\quad (b) $(q,r,p,q)$\quad (c) $(r,q,p,r,q)$\quad (d) $(t,s,p,q,r,p)$\\
(e) $(p,t)$\quad (f) $(p,s,p)$\quad (g) $(u)$
\end{center}
\item List the connected components of $H$. Is $\{p,q,r\}$ a component? Is $H$ connected?
\item Let $D$ be the directed graph with $V=\{1,2,3,4\}$ and $E=\{(1,2),(2,1),(2,3),(4,4),(4,3)\}$. Give the indegree and outdegree of each vertex. What are the tail and head of $(2,3)$? Is $(1,2,3)$ a directed walk? Is $(3,2,1)$?
\item Which of these graphs on $V=\{1,2,3,4,5\}$ are trees?\\
(a) $E=\{\{1,2\},\{2,3\},\{3,4\},\{4,5\}\}$\quad (b) $E=\{\{1,2\},\{2,3\},\{3,1\},\{4,5\}\}$\\
(c) $E=\{\{1,2\},\{1,3\},\{1,4\}\}$\quad (d) $E=\{\{1,2\},\{1,3\},\{1,4\},\{4,5\}\}$
\item Root the tree with $E=\{\{a,b\},\{a,c\},\{b,d\},\{b,e\},\{c,f\},\{f,g\},\{f,h\}\}$ at $a$. List the leaves and the branches. Give depth$(g)$, height$(b)$, height$(c)$ and the depth of the tree. Is it binary? Full?
\item Re-root the tree from problem 7 at $f$. What is the parent of $a$ now? What is the depth of $d$?
\item True or false: (a) every trail is a path; (b) every path is a trail; (c) every cycle is a circuit; (d) every circuit is a cycle; (e) the trivial path $(v)$ is a cycle.
\item In a plain (undirected) graph, explain why no cycle can have exactly two edges. Does the same argument work in a directed graph?
\item Each tree below is binary and rooted at $0$. Which are full? Which are complete?\\
(a) $\{\{0,1\},\{0,2\}\}$\quad (b) $\{\{0,1\},\{0,2\},\{2,3\},\{2,4\}\}$\quad (c) $\{\{0,1\},\{1,2\}\}$\\
(d) $\{\{0,1\},\{0,2\},\{1,3\},\{1,4\},\{2,5\},\{2,6\}\}$
\end{enumerate}

% ------------------------------------------------------------------
\needspace{4\baselineskip}
\subsection{Answers}

\begin{enumerate}
\item (a) Yes. (b) Yes, with one edge: $\{1,2\}$ and $\{2,1\}$ are the same set. (c) No: $3\notin V$, so $\{1,3\}$ isn't a set of vertices of this graph. (d) No: $\{2,2\}=\{2\}$ isn't a two-member set.
\item $\deg(p)=3$ (neighbors $q,r,s$), $\deg(q)=2$, $\deg(r)=2$, $\deg(s)=2$, $\deg(t)=1$, $\deg(u)=0$. The degree of $H$ is the maximum, $3$. Check: $3+2+2+2+1+0=10=2\cdot 5$ edges.
\item \begin{enumerate}
\item Path: edges $\{q,p\},\{p,s\},\{s,t\}$ all exist, no vertex repeats.
\item Trail, not path (the vertex $q$ repeats). It's closed, non-empty, and only first $=$ last repeats, so it's a cycle, and also a circuit and a closed walk.
\item Walk, not trail: the edge $\{r,q\}$ is used first and last (the second time as $\{q,r\}$, which is the same edge).
\item Trail, not path: the five edges are all different, but $p$ appears twice. It isn't closed.
\item Not a walk: there's no edge $\{p,t\}$.
\item Walk, not trail: $\{p,s\}$ is used twice. It's a closed walk, but not a circuit or cycle.
\item The trivial path. It's also a trivial closed walk, but not a circuit or cycle, since it has no edges.
\end{enumerate}
\item The components are $\{p,q,r,s,t\}$ and $\{u\}$. $\{p,q,r\}$ is not a component: it's connected, but you can add $s$ and still be connected, so it isn't maximal. $H$ isn't connected; there's no path from $u$ to anything else.
\item Indegree / outdegree: vertex $1$: $1/1$; vertex $2$: $1/2$; vertex $3$: $2/0$; vertex $4$: $1/2$ (the loop $(4,4)$ counts once each way). Both columns add up to $5$, the number of edges. $(2,3)$ has tail $2$ and head $3$. $(1,2,3)$ is a directed walk, since $(1,2)$ and $(2,3)$ are edges. $(3,2,1)$ isn't: $(3,2)$ is not an edge, and in a directed graph you can't follow $(2,3)$ backwards.
\item (a) Tree: a single line, connected, no cycle. (b) Not a tree: $(1,2,3,1)$ is a cycle, and it's also disconnected. (c) Not a tree: vertex $5$ has no edges, so the graph isn't connected. (d) Tree: connected, and with no cycle.
\item Children: $a\to b,c$; $b\to d,e$; $c\to f$; $f\to g,h$. Leaves: $d,e,g,h$. Branches: $(a,b,d)$, $(a,b,e)$, $(a,c,f,g)$, $(a,c,f,h)$. depth$(g)=3$ (edges $ac$, $cf$, $fg$). height$(b)=1$. height$(c)=2$, via $(c,f,g)$. The depth of the tree is $3$. Binary: yes, no vertex has more than $2$ children. Full: no, $c$ isn't a leaf and has only one child.
\item The path from $a$ to $f$ is $(a,c,f)$, so the parent of $a$ is $c$. The path from $d$ to $f$ is $(d,b,a,c,f)$, four edges, so depth$(d)=4$.
\item (a) False: $(d,c,a,b,c)$ in the running example $G$ is a trail but not a path. (b) True: a repeated edge would force a repeated vertex. (c) True: a cycle is a non-empty trail with first $=$ last, which is exactly what a circuit needs. (d) False: a circuit can revisit a vertex in the middle. (e) False: a cycle must have at least one edge.
\item A two-edge closed walk has the form $(v,w,v)$. Both steps use the edge $\{v,w\}$, so it repeats an edge and isn't a trail, let alone a cycle. In a directed graph the argument breaks: $(v,w)$ and $(w,v)$ are different edges, so $(v,w,v)$ repeats no edge. So a directed graph can have a two-edge cycle, and even a one-edge one from a loop.
\item (a) Full and complete: the root has $2$ children, and both leaves are at depth $1$. (b) Full, not complete: every non-leaf has $2$ children, but the leaves are at depths $1$, $2$ and $2$. (c) Neither: vertices $0$ and $1$ each have just one child. (d) Full and complete: all four leaves are at depth $2$.
\end{enumerate}
